\documentclass[11pt,a4paper]{article}

\usepackage[T1]{fontenc}
\usepackage[utf8]{inputenc}
\usepackage[english]{babel}
\usepackage{times}
\usepackage[margin=25mm]{geometry}
\usepackage{amsmath,amssymb}
\usepackage{siunitx}
\usepackage{booktabs}
\usepackage{longtable}
\usepackage{graphicx}
\usepackage{circuitikz}
\usepackage{caption}
\usepackage{enumitem}
\usepackage[hidelinks]{hyperref}

\sisetup{
  per-mode = symbol,
  range-units = single,
  detect-weight = true
}

\newcommand{\Vces}{V_{\mathrm{CES}}}
\newcommand{\Vdc}{V_{\mathrm{DC}}}
\newcommand{\Vpk}{V_{\mathrm{pk}}}
\newcommand{\Ls}{L_{\sigma}}
\newcommand{\Tvj}{T_{\mathrm{vj}}}

\title{Snubber Design Guidelines for a Half-Bridge Leg\\
Based on the Infineon FZ2400R33HE4 IHM-B Module}
\author{Davide Bagnara}
\date{\today}

\begin{document}
\maketitle

\begin{abstract}
This note collects the design guidelines used to size a snubber network for a
half-bridge leg built with two \SI{3300}{\volt} / \SI{2400}{\ampere}
Trench/Fieldstop IGBT4 modules (Infineon FZ2400R33HE4, IHM-B package).
It defines the voltage budget, quantifies the commutation-loop inductance,
compares the practical snubber topologies, provides closed-form sizing
equations, and works through a numerical example including loss estimation,
component selection and test verification. A final section explains how to
rescale the results to a \SI{1700}{\volt} device of the same current class.
\end{abstract}

\tableofcontents

%======================================================================
\section{Scope and device data}

\subsection{Note on the part number}

The request referred to a \emph{FZ2400R17} device, while the attached datasheet
is the \textbf{FZ2400R33HE4}, i.e.\ a \SI{3300}{\volt} class module.
All numerical results in this note are derived from the attached datasheet
(Revision 1.30, 2022-11-22). Section~\ref{sec:rescale} gives the rescaling
rules for a \SI{1700}{\volt} module of the same current rating, for which the
methodology is identical but the voltage budget and the resulting component
values differ substantially.

\subsection{Data relevant to snubber design}

Only the parameters that enter the snubber calculation are reproduced here.

\begin{table}[htbp]
\centering
\caption{FZ2400R33HE4 parameters relevant to snubber sizing}
\label{tab:device}
\begin{tabular}{llll}
\toprule
Parameter & Symbol & Value & Condition \\
\midrule
Collector--emitter voltage      & $\Vces$        & \SI{3300}{\volt}   & $\Tvj = \SIrange{-40}{150}{\celsius}$ \\
DC stability (cosmic ray)       & $V_{CE(D)}$    & \SI{2100}{\volt}   & $\Tvj = \SI{25}{\celsius}$, 100 FIT \\
Nominal DC collector current    & $I_{C\,nom}$   & \SI{2400}{\ampere} & $T_C = \SI{105}{\celsius}$ \\
Repetitive peak current         & $I_{CRM}$      & \SI{4800}{\ampere} & $t_p$ limited by $\Tvj$ \\
Short-circuit current           & $I_{SC}$       & \SI{9600}{\ampere} & $V_{CC}=\SI{2400}{\volt}$, $t_p\le\SI{10}{\micro\second}$ \\
Module stray inductance         & $L_{sCE}$      & \SI{6}{\nano\henry} & per switch \\
Test circuit stray inductance   & $\Ls$          & \SI{85}{\nano\henry} & datasheet switching test \\
Turn-off fall time              & $t_f$          & \SI{1.4}{\micro\second} & \SI{150}{\celsius}, $R_{G,off}=\SI{3.3}{\ohm}$ \\
Turn-off delay time             & $t_{d,off}$    & \SI{4.3}{\micro\second} & \SI{150}{\celsius} \\
Turn-off $\mathrm{d}v/\mathrm{d}t$        & $\dot v$       & \SI{1500}{\volt\per\micro\second} & $\Tvj = \SI{150}{\celsius}$ \\
Turn-on $\mathrm{d}i/\mathrm{d}t$         & $\dot\imath$   & \SI{7600}{\ampere\per\micro\second} & $\Tvj = \SI{150}{\celsius}$ \\
Internal gate resistance        & $R_{G,int}$    & \SI{0.5}{\ohm} & \\
Input capacitance               & $C_{ies}$      & \SI{280}{\nano\farad} & \SI{25}{\volt}, \SI{1}{\mega\hertz} \\
Reverse transfer capacitance    & $C_{res}$      & \SI{8}{\nano\farad} & \SI{25}{\volt}, \SI{1}{\mega\hertz} \\
Turn-off energy                 & $E_{off}$      & \SI{4950}{\milli\joule} & \SI{2400}{\ampere}, \SI{1800}{\volt}, \SI{150}{\celsius} \\
Diode peak recovery current     & $I_{RM}$       & \SI{2880}{\ampere} & \SI{150}{\celsius}, \SI{7600}{\ampere\per\micro\second} \\
Diode recovered charge          & $Q_r$          & \SI{2500}{\micro\coulomb} & \SI{150}{\celsius} \\
Diode minimum on-time           & $t_{on,min}$   & \SI{10}{\micro\second} & \\
\bottomrule
\end{tabular}
\end{table}

Two entries in Table~\ref{tab:device} deserve emphasis because they dominate
the design:

\begin{itemize}[nosep]
\item The datasheet specifies the short-circuit rating as
      $V_{CE,\max} = \Vces - L_{sCE}\,\mathrm{d}i/\mathrm{d}t$. The
      manufacturer therefore explicitly accounts for the inductive overshoot
      \emph{inside} the module; every nanohenry added by the busbar reduces
      the usable blocking voltage in the same way.
\item The DC stability limit $V_{CE(D)} = \SI{2100}{\volt}$ (100 FIT) caps the
      \emph{steady-state} DC-link voltage. It is a cosmic-ray failure-rate
      limit, unrelated to the transient peak, and it must be checked
      independently.
\end{itemize}

%======================================================================
\section{The overvoltage mechanism}

\subsection{Commutation loop}

At turn-off the collector current is interrupted in the loop formed by the
DC-link capacitor, the busbar, the module terminals and the internal module
leads. The total loop inductance
\begin{equation}
L_{loop} = L_{cap,ESL} + L_{busbar} + 2\,L_{sCE}
\label{eq:loop}
\end{equation}
converts the current slope into a voltage transient across the device:
\begin{equation}
\Vpk = \Vdc + L_{loop}\,\left|\frac{\mathrm{d}i}{\mathrm{d}t}\right|_{\max}
       + \Delta V_{ring}.
\label{eq:vpk}
\end{equation}

\begin{figure}[htbp]
\centering
\begin{circuitikz}[scale=0.95, transform shape]
  % DC link capacitor
  \draw (0,1) to[C, l=$C_{DC}$] (0,5);
  \draw (0,5) -- (1.2,5);
  \draw (0,1) -- (1.2,1);
  % busbar inductance
  \draw (1.2,5) to[L, l=$L_{busbar}/2$] (3.8,5);
  \draw (1.2,1) to[L, l_=$L_{busbar}/2$] (3.8,1);
  % module boundary
  \draw[dashed] (4.3,0.4) rectangle (9.4,5.6);
  \node at (6.85,5.9) {\small module};
  % module internal inductance, upper
  \draw (3.8,5) -- (4.3,5);
  \draw (4.6,5) to[L, l_=$L_{sCE}$] (6.0,5);
  \draw (4.3,5) -- (4.6,5);
  % high-side IGBT
  \draw (6.0,5) to[Tnigbt] (6.0,3.4);
  % high-side diode
  \draw (6.0,5) -- (8.2,5);
  \draw (8.2,3.4) to[D, l_=$D_1$] (8.2,5);
  \draw (6.0,3.4) -- (8.2,3.4);
  % low-side IGBT
  \draw (6.0,3.4) to[Tnigbt] (6.0,1.8);
  % low-side diode
  \draw (6.0,1.8) -- (8.2,1.8);
  \draw (8.2,1.8) to[D, l_=$D_2$] (8.2,3.4);
  % module internal inductance, lower
  \draw (6.0,1.8) -- (6.0,1) to[L, l_=$L_{sCE}$] (4.6,1) -- (3.8,1);
  % AC output
  \draw (8.2,3.4) -- (10.2,3.4) node[right]{\small AC};
  % loop current arrow
  \draw[->,very thick,red] (1.9,4.5) -- (1.9,1.5);
  \node[red,right] at (2.0,3.0) {$i_{loop}$};
\end{circuitikz}
\caption{Commutation loop of the leg. The snubber must be inserted so that it
shortens the loop seen by the switching device, i.e.\ as close as possible to
the module DC terminals.}
\label{fig:loop}
\end{figure}

\subsection{Which current slope applies}

An IGBT does not chop the current instantaneously: the effective slope is set
by the gate drive and by the device physics. Three operating cases have to be
distinguished, and they lead to very different requirements.

\begin{table}[htbp]
\centering
\caption{Current slope estimates for the three dimensioning cases}
\label{tab:didt}
\begin{tabular}{lccc}
\toprule
Case & Current & $|\mathrm{d}i/\mathrm{d}t|$ & $\Vdc$ \\
\midrule
Nominal turn-off       & \SI{2400}{\ampere} & $\approx\SI{2.5}{\kilo\ampere\per\micro\second}$ & \SI{1800}{\volt} \\
Overload turn-off      & \SI{4800}{\ampere} & $\approx\SI{4}{\kilo\ampere\per\micro\second}$   & \SI{1800}{\volt} \\
Short-circuit turn-off & \SI{9600}{\ampere} & $\approx\SI{8}{\kilo\ampere\per\micro\second}$   & \SI{2400}{\volt} \\
\bottomrule
\end{tabular}
\end{table}

For the nominal case the slope follows from the fall time,
$|\mathrm{d}i/\mathrm{d}t| \approx k\,I_C/t_f$ with $k \approx 1.5\ldots 2$
because the current waveform is not linear:
\begin{equation}
\left|\frac{\mathrm{d}i}{\mathrm{d}t}\right| \approx
1.5\cdot\frac{\SI{2400}{\ampere}}{\SI{1.4}{\micro\second}}
\approx \SI{2.5}{\kilo\ampere\per\micro\second}.
\end{equation}

The short-circuit case is far more aggressive: the desaturation protection
turns off a current up to $I_{SC}=\SI{9600}{\ampere}$, and even with a soft
turn-off gate drive the slope reaches several \si{\kilo\ampere\per\micro\second}.

\subsection{Diode recovery}

The second overvoltage source is the reverse recovery of the complementary
diode. With $Q_r = \SI{2500}{\micro\coulomb}$ and
$I_{RM} = \SI{2880}{\ampere}$ at \SI{150}{\celsius}, the recovery current
decay at the end of the recovery produces a transient across the diode.
The emitter-controlled~4 diode is intentionally soft, but snappy behaviour can
still appear at \emph{low current} and \emph{low junction temperature} --- the
worst case for the diode is therefore the opposite of the worst case for the
IGBT. The datasheet minimum on-time $t_{on,min}=\SI{10}{\micro\second}$ must
be respected, otherwise the recovery becomes uncontrolled.

%======================================================================
\section{Voltage budget}

The design starts by allocating the available blocking voltage.
Table~\ref{tab:budget} is the recommended allocation for a \SI{3300}{\volt}
module, following common industrial practice.

\begin{table}[htbp]
\centering
\caption{Voltage budget for the \SI{3300}{\volt} leg}
\label{tab:budget}
\begin{tabular}{lcl}
\toprule
Quantity & Value & Rationale \\
\midrule
$\Vces$                                & \SI{3300}{\volt} & datasheet \\
Nominal DC link $\Vdc$                 & \SI{1800}{\volt} & $\approx 0.55\,\Vces$ \\
Maximum DC link (regulation, braking)  & \SI{2000}{\volt} & below $V_{CE(D)}$ \\
Cosmic-ray limit $V_{CE(D)}$           & \SI{2100}{\volt} & 100 FIT, steady state \\
Design peak limit, normal operation    & \SI{2600}{\volt} & $\approx 0.79\,\Vces$ \\
Design peak limit, fault turn-off      & \SI{3000}{\volt} & $\approx 0.91\,\Vces$, non-repetitive \\
\bottomrule
\end{tabular}
\end{table}

From the peak limits the maximum admissible loop inductance follows directly
by inverting~\eqref{eq:vpk}:
\begin{equation}
L_{loop,\max} = \frac{\Vpk^{lim} - \Vdc}{|\mathrm{d}i/\mathrm{d}t|_{\max}}.
\label{eq:lmax}
\end{equation}

\begin{table}[htbp]
\centering
\caption{Maximum admissible loop inductance per case}
\label{tab:lmax}
\begin{tabular}{lccc}
\toprule
Case & $\Vpk^{lim} - \Vdc$ & $|\mathrm{d}i/\mathrm{d}t|$ & $L_{loop,\max}$ \\
\midrule
Nominal       & \SI{800}{\volt} & \SI{2.5}{\kilo\ampere\per\micro\second} & \SI{320}{\nano\henry} \\
Overload      & \SI{800}{\volt} & \SI{4}{\kilo\ampere\per\micro\second}   & \SI{200}{\nano\henry} \\
Short circuit & \SI{600}{\volt} & \SI{8}{\kilo\ampere\per\micro\second}   & \SI{75}{\nano\henry} \\
\bottomrule
\end{tabular}
\end{table}

\textbf{This is the central result of the analysis.} In normal operation a
well-designed laminated busbar (\SIrange{50}{100}{\nano\henry}) already
satisfies the requirement with a comfortable margin, and \emph{no snubber is
strictly necessary}. The dimensioning case is the \emph{short-circuit
turn-off}, where the admissible loop inductance drops to
\SI{75}{\nano\henry} --- a value that is very hard to reach with a real
mechanical layout. The snubber (or an equivalent measure such as active
clamping) is therefore justified primarily by the fault case and by the
damping of the residual ringing, not by nominal operation.

%======================================================================
\section{Snubber topologies}

\begin{table}[htbp]
\centering
\caption{Comparison of snubber topologies for a high-power leg}
\label{tab:topologies}
\begin{tabular}{p{3.0cm}p{4.2cm}p{4.2cm}p{2.6cm}}
\toprule
Topology & Advantages & Drawbacks & Use \\
\midrule
Pure C (decoupling) &
Lossless, simple, two terminals, bolts directly onto the module &
Undamped LC ringing with $L_{loop}$; no hard clamping &
Standard, always present \\
\addlinespace
Series RC across DC terminals &
Damps the ringing, no diode required &
Continuous loss, $R$ must carry the full transient current &
Common addition \\
\addlinespace
RCD clamp per leg &
Hard clamp of the peak, only active above the clamp level, protects during faults &
Requires a fast \SI{3300}{\volt} diode; bulky &
Recommended here \\
\addlinespace
RCD discharge-suppressing per switch &
Also limits $\mathrm{d}v/\mathrm{d}t$ &
Capacitor becomes very large at \SI{2400}{\ampere}; four components per switch &
Not attractive \\
\addlinespace
Undeland clamp &
Single resistor for the whole leg, handles both IGBT and diode transients &
More components, tuning of the discharge path &
Alternative \\
\addlinespace
Active clamping (gate feedback) &
No power components, self-adapting, excellent for SCSOA &
Increases $E_{off}$, requires a dedicated gate driver &
Preferred modern alternative \\
\bottomrule
\end{tabular}
\end{table}

\begin{figure}[htbp]
\centering
\begin{circuitikz}[scale=0.9, transform shape]
  % (a) C snubber
  \draw (0,3.2) node[above]{(a) C};
  \draw (0,3) -- (1.6,3);
  \draw (0,0) -- (1.6,0);
  \draw (0.8,3) to[C, l=$C_s$] (0.8,0);
  % (b) RC snubber
  \draw (3.4,3.2) node[above]{(b) RC};
  \draw (3.4,3) -- (5.0,3);
  \draw (3.4,0) -- (5.0,0);
  \draw (4.2,3) to[C, l=$C_s$] (4.2,1.7);
  \draw (4.2,1.7) to[R, l=$R_s$] (4.2,0);
  % (c) RCD clamp
  \draw (7.4,3.2) node[above]{(c) RCD clamp};
  \draw (7.0,3) -- (9.8,3);
  \draw (7.0,0) -- (9.8,0);
  \draw (7.6,3) to[D, l_=$D_s$] (7.6,1.7);
  \draw (7.6,1.7) to[C, l_=$C_s$] (7.6,0);
  \draw (9.2,3) to[R, l=$R_s$] (9.2,1.7);
  \draw (7.6,1.7) -- (9.2,1.7);
\end{circuitikz}
\caption{The three practical snubber configurations connected across the DC
terminals of the leg. In (c) the resistor returns the clamped energy to the
DC link.}
\label{fig:topologies}
\end{figure}

%======================================================================
\section{Sizing equations}

\subsection{Pure capacitive snubber}

With a capacitor $C_s$ across the DC terminals, the loop becomes a resonant
$L_{loop}C_s$ circuit with characteristic impedance
\begin{equation}
Z_0 = \sqrt{\frac{L_{loop}}{C_s}}, \qquad
f_0 = \frac{1}{2\pi\sqrt{L_{loop}C_s}}.
\end{equation}
If the device chopped the current instantaneously (worst case), the peak
overshoot would be
\begin{equation}
\Delta V = I_{off}\,Z_0 = I_{off}\sqrt{\frac{L_{loop}}{C_s}}
\quad\Longrightarrow\quad
\boxed{\;C_s = \frac{L_{loop}\,I_{off}^{2}}{\Delta V^{2}}\;}
\label{eq:cs}
\end{equation}
which is equivalent to the energy balance
$\tfrac{1}{2}L_{loop}I_{off}^{2} = \tfrac{1}{2}C_s\Delta V^{2}$.
Equation~\eqref{eq:cs} is conservative because a real IGBT dissipates most of
the loop energy internally (this is precisely what $E_{off}$ accounts for);
it should be used as an upper bound.

Two secondary constraints must be checked:
\begin{align}
\hat\imath_{C_s} &= C_s \left.\frac{\mathrm{d}v}{\mathrm{d}t}\right|_{\max}
  && \text{(peak capacitor current)} \label{eq:icap}\\
f_0 &\ll \frac{1}{2\pi\,t_f}
  && \text{(resonance well above the switching transient)} \nonumber
\end{align}

\subsection{RC damping snubber}

Adding a series resistor turns the loop into a damped second-order system with
damping factor
\begin{equation}
\zeta = \frac{R_s}{2}\sqrt{\frac{C_s}{L_{loop}}} = \frac{R_s}{2 Z_0}.
\end{equation}
Critical damping requires $R_s = 2Z_0$; in practice
$\zeta = 0.5\ldots 1$, i.e.
\begin{equation}
\boxed{\;Z_0 \le R_s \le 2 Z_0\;}
\end{equation}
is a good compromise between ringing suppression and the additional voltage
step $R_s\,I_{off}$ that the resistor itself injects. That step is the reason
why $R_s$ cannot be increased arbitrarily:
\begin{equation}
\Vpk \approx \Vdc + R_s I_{off} + \Delta V_{C}.
\end{equation}

The dissipation depends on how much of the loop energy actually reaches the
resistor. The pessimistic bound and the realistic estimate are
\begin{align}
P_{R,\max} &= \tfrac{1}{2}L_{loop} I_{off}^{2}\, n f_{sw} && \text{(resistor absorbs the whole loop energy)} \label{eq:prmax}\\
P_{R,typ}  &= \tfrac{1}{2}C_s \Delta V_{ring}^{2}\, n f_{sw} && \text{(resistor damps the residual ringing only)} \label{eq:prtyp}
\end{align}
with $n$ the number of damped commutations per switching period ($n=2$ for a
leg: IGBT turn-off and diode recovery).

\subsection{RCD clamp}

The clamp capacitor is pre-charged to $\Vdc$ through $D_s$ and absorbs the
loop energy only when the terminal voltage exceeds $\Vdc$:
\begin{equation}
\tfrac{1}{2}L_{loop} I_{off}^{2}
= \tfrac{1}{2}C_s\left(V_{cl}^{2}-\Vdc^{2}\right)
\quad\Longrightarrow\quad
\boxed{\;C_s = \frac{L_{loop}\,I_{off}^{2}}{V_{cl}^{2}-\Vdc^{2}}\;}
\label{eq:ccl}
\end{equation}
The discharge resistor must return the stored charge to the DC link before the
next commutation:
\begin{equation}
\tau = R_s C_s \le \frac{1}{5 f_{sw}}
\quad\Longrightarrow\quad
\boxed{\;R_s \le \frac{1}{5 f_{sw} C_s}\;}
\label{eq:rcl}
\end{equation}
and its dissipation, when the clamp is active, is
\begin{equation}
P_{R} = \tfrac{1}{2}L_{loop}I_{off}^{2}\,n f_{sw}.
\label{eq:pcl}
\end{equation}
The design intent is to set $V_{cl}$ \emph{above} the normal-operation peak so
that the clamp does not conduct at all in steady state and~\eqref{eq:pcl}
applies only during overload and fault events.

\subsection{Discharge-suppressing RCD (for reference)}

For the per-switch variant the capacitor follows from the requirement that the
collector voltage does not rise faster than the current falls:
\begin{equation}
C_s \ge \frac{I_{off}\,t_f}{2\,\Delta V}.
\end{equation}
At $I_{off}=\SI{2400}{\ampere}$, $t_f=\SI{1.4}{\micro\second}$ and
$\Delta V = \SI{400}{\volt}$ this gives $C_s \ge \SI{4.2}{\micro\farad}$
\emph{per switch}, which confirms that this topology is impractical at this
current level.

%======================================================================
\section{Worked example}

\subsection{Assumptions}

\begin{table}[htbp]
\centering
\caption{Design assumptions}
\label{tab:assump}
\begin{tabular}{lc}
\toprule
Quantity & Value \\
\midrule
Nominal DC link $\Vdc$              & \SI{1800}{\volt} \\
Maximum DC link                      & \SI{2000}{\volt} \\
Switching frequency $f_{sw}$        & \SI{500}{\hertz} \\
Busbar inductance (laminated)        & \SI{80}{\nano\henry} \\
Module contribution $2L_{sCE}$      & \SI{12}{\nano\henry} \\
Capacitor ESL                        & \SI{8}{\nano\henry} \\
\textbf{Total loop inductance $L_{loop}$} & \textbf{\SI{100}{\nano\henry}} \\
\bottomrule
\end{tabular}
\end{table}

\subsection{Unsnubbed peak voltage}

Using~\eqref{eq:vpk} and Table~\ref{tab:didt}:
\begin{align}
\text{Nominal:}\quad
\Vpk &= 1800 + 100\times10^{-9}\cdot 2.5\times10^{9}
      = 1800 + 250 = \SI{2050}{\volt} \\
\text{Overload:}\quad
\Vpk &= 1800 + 100\times10^{-9}\cdot 4\times10^{9}
      = 1800 + 400 = \SI{2200}{\volt} \\
\text{Short circuit:}\quad
\Vpk &= 2400 + 100\times10^{-9}\cdot 8\times10^{9}
      = 2400 + 800 = \SI{3200}{\volt}
\end{align}

Normal and overload operation stay well below the \SI{2600}{\volt} design
limit. The short-circuit turn-off exceeds the \SI{3000}{\volt} fault limit and
approaches $\Vces$: it requires either a clamp or an active-clamping gate
driver.

\subsection{Decoupling capacitor}

Selecting a decoupling film capacitor $C_s = \SI{1}{\micro\farad}$:
\begin{align}
Z_0 &= \sqrt{\frac{\SI{100}{\nano\henry}}{\SI{1}{\micro\farad}}}
     = \SI{0.316}{\ohm} \\
f_0 &= \frac{1}{2\pi\sqrt{100\times10^{-9}\cdot 1\times10^{-6}}}
     = \SI{503}{\kilo\hertz}
\end{align}
The resonant frequency is two decades above the reciprocal of the fall time
($1/2\pi t_f \approx \SI{114}{\kilo\hertz}$), which is the correct
separation. The peak current through the capacitor,
from~\eqref{eq:icap} with the datasheet
$\mathrm{d}v/\mathrm{d}t = \SI{1500}{\volt\per\micro\second}$, is
\begin{equation}
\hat\imath_{C_s} = \SI{1}{\micro\farad}\cdot\SI{1500}{\volt\per\micro\second}
= \SI{1500}{\ampere},
\end{equation}
which is a hard constraint on the capacitor selection and the reason why a
single high-$\mathrm{d}v/\mathrm{d}t$ film capacitor, or several in parallel,
is required.

\subsection{Damping resistor}

\begin{equation}
R_s = Z_0 \ldots 2Z_0 = \SIrange{0.32}{0.63}{\ohm}
\quad\Rightarrow\quad R_s = \SI{0.5}{\ohm}\ (\zeta = 0.79)
\end{equation}
The voltage step introduced by the resistor at nominal current is
$R_s I_{off} = 0.5\times 2400 = \SI{1200}{\volt}$, which is unacceptable if
the resistor carries the full current. This is a classical trap: the
\emph{series} RC snubber must be sized with a much smaller $R_s$, or the
resistor must be placed only in the damping path of an RCD structure. A
practical compromise for a pure RC across the DC terminals is
$R_s = \SIrange{0.1}{0.2}{\ohm}$ with $C_s = \SI{2}{\micro\farad}$, accepting
$\zeta \approx 0.35$ and a residual ringing of a few oscillations.

Dissipation from~\eqref{eq:prtyp} with $\Delta V_{ring} = \SI{250}{\volt}$,
$n=2$, $f_{sw}=\SI{500}{\hertz}$:
\begin{equation}
P_{R,typ} = \tfrac{1}{2}\cdot 2\times10^{-6}\cdot 250^{2}\cdot 2\cdot 500
= \SI{62}{\watt}.
\end{equation}
The pessimistic bound~\eqref{eq:prmax} gives
$\tfrac12\cdot100\times10^{-9}\cdot2400^{2}\cdot2\cdot500 = \SI{288}{\watt}$,
which sets the thermal rating of the resistor assembly.

\subsection{RCD clamp for the fault case}

Clamping level $V_{cl}=\SI{2800}{\volt}$, $\Vdc = \SI{2400}{\volt}$ (fault
condition), $I_{off}=I_{SC}=\SI{9600}{\ampere}$, from~\eqref{eq:ccl}:
\begin{equation}
C_s = \frac{100\times10^{-9}\cdot 9600^{2}}{2800^{2}-2400^{2}}
    = \frac{9.216}{2.08\times10^{6}}
    = \SI{4.4}{\micro\farad}
\end{equation}
Choose $C_s = \SI{5}{\micro\farad}$ (film, \SI{3000}{\volt} DC rating).
Discharge resistor from~\eqref{eq:rcl}:
\begin{equation}
R_s \le \frac{1}{5\cdot 500\cdot 5\times10^{-6}} = \SI{80}{\ohm}
\quad\Rightarrow\quad R_s = \SI{68}{\ohm}
\end{equation}
Since $V_{cl}=\SI{2800}{\volt}$ is above the \SI{2200}{\volt} overload peak,
the clamp is idle in normal operation and the resistor only has to absorb the
single-event energy
\begin{equation}
E_{ev} = \tfrac{1}{2} L_{loop} I_{SC}^{2}
= \tfrac{1}{2}\cdot 100\times10^{-9}\cdot 9600^{2} = \SI{4.6}{\joule},
\end{equation}
plus the energy already stored in $C_s$. A pulse-rated wirewound or thick-film
resistor of a few tens of joules single-pulse capability is sufficient.

\subsection{Summary of the design}

\begin{table}[htbp]
\centering
\caption{Resulting snubber design for the leg}
\label{tab:result}
\begin{tabular}{llll}
\toprule
Function & Component & Value & Rating \\
\midrule
Decoupling / damping & $C_s$ & \SIrange{1}{2}{\micro\farad} & \SI{2500}{\volt} DC, ESL $<\SI{20}{\nano\henry}$, $\hat\imath > \SI{2}{\kilo\ampere}$ \\
Damping              & $R_s$ & \SIrange{0.1}{0.2}{\ohm}     & non-inductive, \SI{300}{\watt} \\
Clamp capacitor      & $C_{cl}$ & \SI{5}{\micro\farad}      & \SI{3000}{\volt} DC film \\
Clamp diode          & $D_{cl}$ & --                        & \SI{4500}{\volt}, fast, soft recovery \\
Clamp resistor       & $R_{cl}$ & \SI{68}{\ohm}             & \SI{50}{\joule} single pulse \\
\bottomrule
\end{tabular}
\end{table}

%======================================================================
\section{Component selection}

\subsection{Capacitor}

\begin{itemize}[nosep]
\item Dry metallized polypropylene film, self-healing. Ceramic is unsuitable
      at this energy and voltage.
\item DC voltage rating $\ge 1.3\ldots1.5\times$ the maximum applied DC
      voltage.
\item $\mathrm{d}v/\mathrm{d}t$ rating: verify against~\eqref{eq:icap}. At
      \SI{1500}{\volt\per\micro\second} a \SI{1}{\micro\farad} capacitor sees
      \SI{1500}{\ampere} peak.
\item ESL: the capacitor and its connection must contribute a small fraction
      of $L_{loop}$. Bolt-on terminals matching the module M8 studs, not
      leaded packages.
\item RMS current: for the pure decoupling capacitor the RMS current is a few
      tens of amperes; verify against the ripple spectrum.
\end{itemize}

\subsection{Resistor}

\begin{itemize}[nosep]
\item Non-inductive construction (thick film on ceramic, carbon composition,
      or bifilar wirewound). Series inductance $<\SI{50}{\nano\henry}$,
      otherwise the resistor defeats its own purpose.
\item Single-pulse energy rating from the fault case, average power rating
      from~\eqref{eq:prmax}.
\item Several resistors in parallel, symmetrically arranged, help both the
      thermal rating and the inductance.
\end{itemize}

\subsection{Clamp diode}

\begin{itemize}[nosep]
\item $V_{RRM} \ge \SI{4500}{\volt}$ to keep margin over the clamp level.
\item Fast and \emph{soft} recovery: a snappy clamp diode generates its own
      overvoltage and can be worse than no clamp.
\item Low forward recovery voltage $V_{FR}$: the diode must not overshoot
      during turn-on, otherwise the clamp is ineffective in the first
      microsecond, which is exactly when it is needed.
\end{itemize}

%======================================================================
\section{Layout guidelines}

\begin{enumerate}[nosep]
\item Laminated busbar with maximum overlap and minimum dielectric thickness
      compatible with the insulation coordination
      (Table~1 of the datasheet: creepage \SI{32.2}{\milli\metre}, clearance
      \SI{19.1}{\milli\metre} terminal-to-heatsink).
\item Snubber capacitor bolted directly to the module DC terminals, not to the
      DC-link capacitor. The purpose is to shorten the loop seen by the chip.
\item Symmetric arrangement if modules are paralleled: an asymmetry in
      $L_{loop}$ translates into unequal turn-off overvoltage and unequal
      current sharing.
\item Keep the gate-drive return path separate from the power loop; the
      auxiliary emitter terminals of the IHM-B package exist for this reason.
\item Verify the loop inductance experimentally rather than trusting a
      calculation: from a double-pulse test,
      $L_{loop} = \Delta V / (\mathrm{d}i/\mathrm{d}t)$ measured at the
      terminals.
\end{enumerate}

%======================================================================
\section{Alternatives to a passive snubber}

Given the results of Table~\ref{tab:lmax}, two gate-side measures may replace
or complement the snubber:

\begin{description}[nosep,style=nextline]
\item[Active clamping] A chain of transient-voltage suppressors between
      collector and gate feeds current back into the gate when $V_{CE}$
      exceeds the clamp level, slowing the turn-off. It is compact, requires
      no bulky passives, and is very effective for the short-circuit case. The
      penalty is an increase of $E_{off}$ and of the device junction
      temperature during the clamped event.
\item[Two-level (soft) turn-off] The driver reduces the gate voltage to an
      intermediate level before switching off completely, or increases
      $R_{G,off}$ during a detected fault. With $R_{G,off}$ increased from
      \SI{3.3}{\ohm} to \SI{10}{\ohm} the fault $\mathrm{d}i/\mathrm{d}t$ can
      typically be halved, which relaxes~\eqref{eq:lmax} by the same factor.
      The datasheet switching-loss and switching-time curves
      $E = f(R_G)$ and $t = f(R_G)$ quantify the trade-off.
\end{description}

Increasing $R_{G,off}$ is the cheapest lever available and should be evaluated
before adding copper and components: it costs switching energy but no volume.

%======================================================================
\section{Verification plan}

\begin{enumerate}[nosep]
\item \textbf{Loop inductance measurement.} Double-pulse test at low voltage;
      extract $L_{loop}$ from $\Delta V/(\mathrm{d}i/\mathrm{d}t)$ and compare
      with the FEM/analytical busbar estimate.
\item \textbf{Nominal double-pulse test.} $\Vdc = \SI{2000}{\volt}$,
      $I_C = \SI{2400}{\ampere}$, $\Tvj = \SI{150}{\celsius}$. Verify
      $\Vpk \le \SI{2600}{\volt}$ and check the waveform against the
      datasheet $t_{d,off}$, $t_f$ and $E_{off}$.
\item \textbf{Overload test.} $I_C = I_{CRM} = \SI{4800}{\ampere}$; verify the
      trajectory stays inside the RBSOA diagram of the datasheet.
\item \textbf{Diode recovery test.} Low current (\SIrange{10}{20}{\percent} of
      $I_{nom}$) and low temperature (\SI{25}{\celsius}) to expose snappy
      recovery; verify $t_{on,min} \ge \SI{10}{\micro\second}$ is respected by
      the modulator.
\item \textbf{Short-circuit tests.} Type~I and Type~II at
      $\Vdc = \SI{2400}{\volt}$, $t_p \le \SI{10}{\micro\second}$; verify
      $\Vpk$ against $\Vces - L_{sCE}\,\mathrm{d}i/\mathrm{d}t$ and check the
      clamp energy.
\item \textbf{Thermal check of the snubber.} Measure the resistor temperature
      in continuous operation at maximum modulation index and compare with the
      estimate from~\eqref{eq:prtyp}.
\end{enumerate}

%======================================================================
\section{Rescaling to a \SI{1700}{\volt} module}
\label{sec:rescale}

If the intended device is a \SI{1700}{\volt} / \SI{2400}{\ampere} module, the
methodology is unchanged but the numbers shift as follows.

\begin{table}[htbp]
\centering
\caption{Voltage budget comparison}
\label{tab:rescale}
\begin{tabular}{lcc}
\toprule
Quantity & \SI{3300}{\volt} class & \SI{1700}{\volt} class \\
\midrule
$\Vces$                             & \SI{3300}{\volt} & \SI{1700}{\volt} \\
Typical DC link                     & \SI{1800}{\volt} & \SI{900}{\volt} \\
Design peak, normal operation       & \SI{2600}{\volt} & \SI{1350}{\volt} \\
Available headroom $\Vpk - \Vdc$   & \SI{800}{\volt}  & \SI{450}{\volt} \\
Typical $t_f$                       & $\approx\SI{1.4}{\micro\second}$ & $\approx\SI{0.4}{\micro\second}$ \\
\bottomrule
\end{tabular}
\end{table}

Two effects act in the same, unfavourable direction:
\begin{itemize}[nosep]
\item the available headroom is roughly halved;
\item the \SI{1700}{\volt} device is intrinsically faster, so
      $\mathrm{d}i/\mathrm{d}t$ at the same current is two to three times
      higher.
\end{itemize}
From~\eqref{eq:lmax} the admissible loop inductance therefore drops by roughly
a factor of five, into the \SIrange{20}{40}{\nano\henry} range for nominal
operation. For a \SI{1700}{\volt} leg at \SI{2400}{\ampere} the busbar design
becomes the dominant engineering effort, and the decoupling capacitor is
mandatory rather than optional. The clamp capacitor
from~\eqref{eq:ccl} scales as $1/(V_{cl}^{2}-\Vdc^{2})$ and grows
substantially: a recalculation with the actual datasheet values of the chosen
part is required.

%======================================================================
\section{Conclusions}

\begin{enumerate}[nosep]
\item For the FZ2400R33HE4 at $\Vdc = \SI{1800}{\volt}$ with a
      \SI{100}{\nano\henry} commutation loop, the nominal turn-off overvoltage
      is about \SI{250}{\volt}, giving a peak of \SI{2050}{\volt} against a
      \SI{2600}{\volt} design limit. A snubber is \emph{not} required for
      overvoltage limitation in normal operation.
\item A decoupling film capacitor of \SIrange{1}{2}{\micro\farad} directly on
      the module terminals is nevertheless mandatory, to shorten the
      commutation loop and to keep the resonance well separated from the
      switching transient.
\item The dimensioning case is the short-circuit turn-off, which would require
      $L_{loop} \le \SI{75}{\nano\henry}$. This is addressed by an RCD clamp
      set at \SI{2800}{\volt}, by active clamping, or by a soft turn-off gate
      drive --- the last being the cheapest option.
\item The series resistance of an RC snubber is constrained by the voltage
      step $R_s I_{off}$; at \SI{2400}{\ampere} this forces
      $R_s \le \SI{0.2}{\ohm}$ and accepts partial damping.
\item All results must be confirmed by a double-pulse test at maximum voltage,
      maximum current and \SI{150}{\celsius}, plus a short-circuit test at
      \SI{2400}{\volt}.
\end{enumerate}

%======================================================================
\begin{thebibliography}{9}

\bibitem{ds}
Infineon Technologies AG,
\emph{FZ2400R33HE4 IHM-B Module Datasheet},
Rev.~1.30, 2022-11-22.

\bibitem{undeland}
T.~M. Undeland, F.~Jenset, A.~Steinbakk, T.~Rogne, and M.~Hernes,
``A snubber configuration for both power transistors and GTO PWM inverters,''
in \emph{Proc. IEEE Power Electronics Specialists Conf. (PESC)},
1984, pp.~42--53.

\bibitem{mohan}
N.~Mohan, T.~M. Undeland, and W.~P. Robbins,
\emph{Power Electronics: Converters, Applications and Design},
3rd ed. Hoboken, NJ, USA: Wiley, 2003, ch.~27.

\bibitem{wintrich}
A.~Wintrich, U.~Nicolai, W.~Tursky, and T.~Reimann,
\emph{Application Manual Power Semiconductors},
2nd ed. Ilmenau, Germany: SEMIKRON International, 2015.

\bibitem{lutz}
J.~Lutz, H.~Schlangenotto, U.~Scheuermann, and R.~De~Doncker,
\emph{Semiconductor Power Devices: Physics, Characteristics, Reliability},
2nd ed. Cham, Switzerland: Springer, 2018.

\bibitem{an}
Infineon Technologies AG,
\emph{AN2008-03: Thermal equivalent circuit models} and
\emph{AN-1103: Mounting instructions for IHM/IHV modules},
Application Notes.

\bibitem{cosmic}
Infineon Technologies AG,
\emph{AN2019-05: Cosmic ray induced failures in power semiconductors},
Application Note.

\end{thebibliography}

\end{document}
